\documentclass[11pt]{article}
\usepackage{amsmath,amssymb,amsthm,mathtools}
\usepackage[margin=1in]{geometry}
\usepackage{booktabs}
\usepackage{hyperref}

\newtheorem{theorem}{Theorem}
\newtheorem{lemma}{Lemma}
\newtheorem{definition}{Definition}
\newtheorem{proposition}{Proposition}
\newtheorem{warning}{Warning}
\newtheorem{obligation}{Obligation}

\newcommand{\C}{\mathcal C}
\newcommand{\B}{\mathcal B}
\newcommand{\Hc}{\mathcal H}
\newcommand{\eps}{\varepsilon}
\newcommand{\sfproj}{\Pi_{\rm sf}}
\newcommand{\XiG}{\Xi_{\rm GCD}}

\title{Step 131: Residual Coefficient Class versus the GCD--Log Blind Sector}
\author{Six Birds RH Membrane Program}
\date{2026-05-14}

\begin{document}
\maketitle

\section{Purpose}
Step 130 showed that the Bui--Pratt--Robles--Zaharescu shifted second-moment main term cannot be imported as an unrestricted matrix lower frame.  The obstruction is an explicit near-null sector of the GCD--log kernel on squarefree Boolean cubes.  Step 131 asks the next, sharper question:
\[
    \sfproj R_NY_{R,N}\stackrel{?}{=}0,
\]
or at least whether this projection is tail-small on the actual residual coefficient class arising from the regularized Burnol/Muentz shadow.

The key point is that the unrestricted coefficient space contains near-null Boolean sign modes, but the actual residual coefficient image may be special.  In particular, the zeta/Muentz shadow has a convolutional structure that can suppress deep Boolean blind modes.  This step isolates the exact residual:
\[
    \boxed{\XiG_N := R_N^*\sfproj R_N.}
\]

\section{Squarefree Boolean cube and GCD--log diagonalization}
Let $\mathcal P_y=\{p_1,\dots,p_k\}$ be a finite set of small primes and let
\[
    \C_y=\{d: d\text{ squarefree and all }p\mid d\text{ lie in }\mathcal P_y\}.
\]
Identifying $d\in\C_y$ with a subset of $\mathcal P_y$, define the Boolean sign basis
\[
    \psi_\epsilon(d)=2^{-k/2}\prod_{p\mid d}\epsilon_p,
    \qquad \epsilon_p\in\{\pm1\}.
\]
On $\ell^2(\C_y)$, consider the normalized GCD--log kernel
\[
    \widetilde K_q(m,n)
    =\frac{(m,n)}{\sqrt{mn}}
    \left(C_q-\log\frac{mn}{(m,n)^2}\right),
    \qquad C_q=\log q+O(1).
\]

\begin{theorem}[Boolean diagonalization]
The Boolean sign basis diagonalizes $\widetilde K_q$.  The eigenvalue corresponding to $\epsilon=(\epsilon_p)_{p\in\mathcal P_y}$ is
\[
    \kappa_\epsilon(q,y)
    =\Lambda_\epsilon(y)
    \left(
      C_q+
      \sum_{p\in\mathcal P_y}
      \frac{-\epsilon_p p^{-1/2}\log p}{1+\epsilon_p p^{-1/2}}
    \right),
\]
where
\[
    \Lambda_\epsilon(y)=\prod_{p\in\mathcal P_y}(1+\epsilon_pp^{-1/2}).
\]
In particular, the all-minus mode has
\[
    \kappa_-(q,y)=
    \prod_{p\le y}(1-p^{-1/2})
    \left(
      C_q+
      \sum_{p\le y}\frac{p^{-1/2}\log p}{1-p^{-1/2}}
    \right),
\]
which can be much smaller than $\log q$ when $y$ grows like a small multiple of $\log q$.
\end{theorem}

\begin{definition}[GCD--log blind sector]
Fix a threshold $\theta>0$ or a bottom spectral quantile.  The finite GCD--log blind sector is
\[
    \Hc_{{\rm sf},y}(\theta)
    =\operatorname{span}\{\psi_\epsilon:\kappa_\epsilon(q,y)\le \theta\log q\}.
\]
Let $\sfproj_{q,y,\theta}$ be the orthogonal projection onto this sector.
\end{definition}

\begin{definition}[Residual GCD adequacy residual]
For a residual coefficient map $R_N:Y_{R,N}\to\ell^2(\C_y)$, define
\[
    \boxed{\XiG_N=R_N^*\sfproj_{q,y,\theta}R_N.}
\]
This is the part of the residual coefficient class that lies in the GCD--log blind sector.
\end{definition}

\begin{proposition}[Restricted import criterion]
The BPRZ shifted-main-term lower-frame import may be applied to the residual coefficient class only after one of the following records is earned:
\[
    \XiG_N=0,
\]
\[
    \XiG_N\preceq E_N,
    \qquad E_N\to0\quad\text{in a fixed/exhaustive ledger sense},
\]
or
\[
    \XiG_N\preceq F_N+E_N
\]
for a declared Hecke/Dirichlet source frame $F_N$ and vanishing defect.  Without such a record, the unrestricted GCD--log near-null sector blocks the lower-frame import.
\end{proposition}

\section{Zeta convolution and Boolean Möbius annihilation}
The regularized Burnol/Muentz shadow is not an arbitrary coefficient vector.  Its zeta factor induces a convolutional shape.  The finite Boolean model of this effect is the zeta-incidence operator
\[
    (Z_yb)(d)=\sum_{e\mid d}b(e),
    \qquad d\in\C_y.
\]
This map has a useful exact relation to Boolean blind modes.

\begin{lemma}[Möbius annihilation identity]
For each sign vector $\epsilon\in\{\pm1\}^{\mathcal P_y}$,
\[
    \langle \psi_\epsilon,Z_yb\rangle
    =2^{-k/2}\left(\bigotimes_{p\in\mathcal P_y}\ell_{\epsilon_p}\right)b,
\]
where the local functionals on a one-prime seed $(b_0,b_1)$ are
\[
    \ell_+(b_0,b_1)=2b_0+b_1,
    \qquad
    \ell_-(b_0,b_1)=-b_1.
\]
In particular, for the all-minus mode,
\[
    \boxed{
    \langle\psi_-,Z_yb\rangle
    =2^{-k/2}(-1)^k b(P_y),
    }
\]
where $P_y=\prod_{p\in\mathcal P_y}p$.
\end{lemma}

\begin{proof}
The operator $Z_y$ is the tensor product over primes of the local incidence map
\[
    Z_p(b_0,b_1)=(b_0,b_0+b_1).
\]
Pairing with the local plus and minus characters gives
\[
    (1,1)Z_p(b_0,b_1)=2b_0+b_1,
\]
\[
    (1,-1)Z_p(b_0,b_1)=-b_1.
\]
Tensoring over $p\in\mathcal P_y$ gives the result.
\end{proof}

\begin{obligation}[High-divisibility / sieve-tail suppression]
Let $b_N$ be the seed coefficient vector arising from the regularized Burnol/Muentz shadow.  The required blind-sector suppression record is an estimate of the form
\[
    \sum_{\epsilon\in\operatorname{Blind}(q,y,\theta)}
    \left|\left(\bigotimes_p\ell_{\epsilon_p}\right)b_N\right|^2
    \le \delta_N\,\|Z_yb_N\|^2,
\]
with $\delta_N\to0$ or at least $\delta_N\le\delta<1$ strong enough to retain a positive visibility floor.
\end{obligation}

This is the new analytic gate.  It is a high-divisibility statement: blind modes with many minus signs only see seed coefficients divisible by many of the small primes in the squarefree cube.  If the Burnol/Muentz seed suppresses such coefficients, the actual residual class may avoid the GCD--log blind sector even though the unrestricted coefficient space does not.

\section{Consequences for the RH source route}
The residual source route after Step 130 was blocked in the unrestricted coefficient space.  Step 131 refines the diagnosis:
\[
    \boxed{
    \text{unrestricted BPRZ lower frame fails, but the actual residual image may avoid the blind sector.}
    }
\]
The route now has three lawful exits:
\begin{enumerate}
    \item Prove $\XiG_N=0$ or $\XiG_N\to0$ by high-divisibility/Möbius suppression for the actual Burnol/Muentz residual coefficients.
    \item Prove a positive floor $\XiG_N\preceq \delta G_{R,N}$ with $\delta<1$, and compensate with a diverging source strength.
    \item If the blind residual persists, add a separate non-smuggled source frame that charges $\XiG_N$.
\end{enumerate}

\section{Finite sanity checks}
The attached finite checks are algebraic sanity checks only.  They verify:
\begin{itemize}
    \item Boolean diagonalization of the GCD--log kernel;
    \item exact Möbius annihilation identity for the all-minus mode;
    \item strong suppression of blind-sector overlap for synthetic zeta-convolved seeds with high-divisibility suppression;
    \item the finite residual form $\XiG_N=R_N^*\sfproj R_N$.
\end{itemize}
They are not RH evidence.

\section{Next gate}
The next mathematical target is a sieve-tail theorem for the actual regularized Burnol/Muentz seed coefficients:
\[
    \boxed{
    \text{prove that Burnol/Muentz residual seeds have small high-divisibility mass on Boolean blind modes.}
    }
\]
If this succeeds, the BPRZ matrix import can be attempted on the restricted residual class.  If it fails, $\Xi_{\rm GCD}$ is the next residual requiring source absorption.

\end{document}
